\section{Fazit}
\label{sec:fazit}

Die vorliegende Arbeit untersuchte, ob aus dem täglichen Sentiment auf X\,(Twitter) ein statistisch signifikanter Zusammenhang mit den Tagesrenditen von zwölf Fokus-Unternehmen aus den Bereichen Energie, Wärmepumpen und Elektromobilität abzuleiten ist. Dafür wurden ein selbst erhobener Tweet-Datensatz und historische Börsenkursdaten mithilfe eines multilingualen Transformer-Modells, eines lexikonbasierten Baseline-Verfahrens sowie vier statistischer Analysemethoden ausgewertet.

Die zentralen Befunde lassen sich wie folgt zusammenfassen: Im beobachteten Zeitraum von rund vier Monaten (Januar bis April 2026) fand sich kein konsistentes Muster, das auf eine systematische Vorhersagbarkeit von Kursrenditen durch Twitter-Sentiment hinweist. Nach Anwendung der Benjamini-Hochberg-FDR-Korrektur für multiples Testen verblieb keine Pearson-Korrelation unter der Signifikanzschwelle von $\alpha = 0{,}05$ -- der einzige rohe Befund (Daikin, Lag~1, $p_{\mathrm{roh}} = 0{,}048$) wies nach Korrektur $p_{\mathrm{adj}} = 0{,}830$ auf und ist als zufälliger Treffer bei 40 simultanen Tests einzuordnen. Granger-Tests blieben für alle Ticker ohne signifikantes Ergebnis. Rolling-Korrelationen zeigten keine stabilen Zusammenhangsphasen. Event-Study-Profile variierten stark zwischen Unternehmen und Richtungen, ohne ein unternehmensübergreifendes Signal zu erzeugen.

Diese Ergebnisse sind in erster Linie durch den kurzen Beobachtungszeitraum und die damit verbundene geringe statistische Power zu erklären. Eine Replikation mit einem längeren Datensatz -- idealerweise über mehrere Jahre und über verschiedene Marktphasen -- wäre notwendig, um belastbarere Aussagen treffen zu können. Darüber hinaus könnte eine höhere zeitliche Auflösung (intraday), eine feingranularere Segmentierung nach Account-Typen sowie eine gezieltere Vorverarbeitung von mehrsprachigen Inhalten die Aussagekraft zukünftiger Analysen verbessern.

Methodisch leistet die Arbeit einen Beitrag, indem sie ein reproduzierbares, modulares Analyse-Pipeline-Design zeigt, das Datenerhebung, Sentiment-Quantifizierung und statistische Modellierung transparent trennt und für jeden Analyseschritt Qualitäts- und Coverage-Prüfungen dokumentiert. Die Implementierung der Retweet-Separation, der Bot-Filterung sowie des Overlap-Reports sind Bestandteile, die sich in ähnlichen Projekten wiederverwenden lassen.
